% =============================================================================
% Chapter 8  Numerical Study  --  ~10.5 pp
% rubric: Project Body (Modelling) -- Results & Discussion /30.
% Plan: ThesisPlanning.md §5.8.
%
% PURPOSE     Evidence every claim of the Abstract and of §1.3, "without
%             reference to any other documents", one section per question or
%             contribution.
% CRITERION   "Complete, clearly presented results, with detailed discussion
%             showing insight... or warranting publication." Every figure is
%             interpreted in prose; each section ends with what it established.
% SOURCES     P8's records (results-v1) through \res; the reports
%             pihm/docs/10_reproduction.md, 11_pillar1.md, 12_pillar2.md,
%             13_descriptor.md, FILTERS.md, EXTENSIONS.md, NAVIER_STOKES.md,
%             18_claims.md; the notebooks of Planning.md §10.8, one per figure.
% WRITE WHEN  Section by section as the gates pass: §8.2 at G1/G3 (W2);
%             §8.3-8.6 at G4 (W3); §8.7 at G5 (W4).
% EVERY NUMBER is a \res from 0_results/ (\prov{...} only while pihm does not
% produce it yet), and EVERY number carries its tier (T1, T2 or T3) in the
% table or caption.
%
% v2.0 MIGRATION (ThesisPlanning.md §4.3 step 4): this chapter is regrouped
% from nine subsections into §5.8's seven. "Experimental Design and Cost
% Conventions" (the old sec:results-design) has MOVED to
% 7_pipeline.tex's new sec:pipeline-design (§4.3 step 2) -- it is a model-
% input decision, not a result. "Operator Verification and Encoding Cost"
% (sec:results-operators) and "State Preparation and Readout"
% (sec:results-stateprep) merge into the new sec:results-faithful. "Scaling"
% keeps its label (sec:results-scaling, cited from Chs 4-5 and the Abstract)
% but becomes §8.3. "Representative Problems" (sec:results-cases) folds into
% the new sec:results-gap, which is §5.5's challenge made a results section
% (§8.4). "Nonlinear Benchmarks" and "Case Study: Navier--Stokes" keep their
% labels; the single figure they used to share (fig:carleman-ns) splits into
% one per section (F13, F14), per §5.8 points 5-6.
% =============================================================================

% Opening: no heading; follows the section heading directly (no stacked headings).
% ~0.25 pp -- BEFORE §8.1. LAND: what each section evidences; the tiers; the
%   records' provenance in one sentence; the full per-problem results ->
%   Appendix J (app:catalogue). No label: this is unnumbered lead-in prose,
%   not a subsection the outline points to.

\subsection{Reproducing Wu et al.}
\label{sec:results-reproduction}
% ~1.75 pp (S1). Framed as REPRODUCTION, never as rivalry (Style Trap 14).
% LAND, in order:
%   1. tab:eta-reproduction: eta_e for the 11 ODE panels (printed, ideal
%      ||b||^2, ours at the printed n) and a summary line for the five others
%      (Appendix K has all 16). The closed form eq:eta-e and what it
%      validates: our conventions are the paper's.
%   2. fig:faithful-vs-corrected (nb 11): the panels as printed against the
%      problem for Fig. 5, the distortion from rounded points and Maclaurin
%      sources.
%   3. The Fig. 6 and 7 NON-REPRODUCTIONS and the hypotheses tested; the
%      Laplace typo; the heat panel's backward data.
%   4. The paper's quantum claims (nb 12): the time rule falls short of
%      epsilon (A-11), the linear H is semi-definite where the discrete
%      solution is exact (A-09), the doubled-space kernel count (A-10).
%   5. The claims verdicts -> Appendix K (app:reproduction; one paragraph
%      here).
% TAB: tab:eta-reproduction (T6). FIG: fig:faithful-vs-corrected (F6).

\begin{equation}
    \eta_e = \|b\|^2 = \frac{N}{\pi}\int_{-1}^{1}\frac{f(x)^2}{\sqrt{1-x^2}}\,dx .
    \label{eq:eta-e}
\end{equation}

\begin{table}[htbp]
    \centering
    \caption{The paper's printed scale factors against the exact solution's $\|b\|^2$ (\Cref{eq:eta-e}) and the paper-faithful ground state at the printed $n$. $^\dagger$: the ground state is not trusted (a degenerate kernel, or a gap ratio below 100). Classical computation.}
    \label{tab:eta-reproduction}
    \small
    \begin{tabular}{@{}l l r l l l l@{}}
        \toprule
        Fig. & Problem & $n$ & Printed $\eta_e$ & Ideal $\|b\|^2$ & Ground state & Agreement \\
        \midrule
        3b & R2a & 3 & 1.29 & \res[places=5]{build/R2a/faithful/standard/n3}{eta_ideal} & $\res[places=5]{build/R2a/faithful/standard/n3}{eta}\resflag{build/R2a/faithful/standard/n3}{trusted}{}{^{\dagger}}$ & \checkmark \\
        3c & R2b & 3 & 32.47 & \res[places=5]{build/R2b/faithful/standard/n3}{eta_ideal} & $\res[places=5]{build/R2b/faithful/standard/n3}{eta}\resflag{build/R2b/faithful/standard/n3}{trusted}{}{^{\dagger}}$ & \checkmark (one unit low) \\
        3d & R2c & 5 & 451.71 & \res[places=5]{build/R2c/faithful/standard/n5}{eta_ideal} & $\res[places=5]{build/R2c/faithful/standard/n5}{eta}\resflag{build/R2c/faithful/standard/n5}{trusted}{}{^{\dagger}}$ & \checkmark \\
        4a & R3a & 3 & 1.75 & \res[places=5]{build/R3a/faithful/standard/n3}{eta_ideal} & $\res[places=5]{build/R3a/faithful/standard/n3}{eta}\resflag{build/R3a/faithful/standard/n3}{trusted}{}{^{\dagger}}$ & \checkmark \\
        4b & R3b & 3 & 1.49 & \res[places=5]{build/R3b/faithful/standard/n3}{eta_ideal} & $\res[places=5]{build/R3b/faithful/standard/n3}{eta}\resflag{build/R3b/faithful/standard/n3}{trusted}{}{^{\dagger}}$ & \checkmark \\
        4c & R3c & 5 & 108.51 & \res[places=5]{build/R3c/faithful/standard/n5}{eta_ideal} & $\res[places=5]{build/R3c/faithful/standard/n5}{eta}\resflag{build/R3c/faithful/standard/n5}{trusted}{}{^{\dagger}}$ & \checkmark \\
        4d & R3d & 5 & 135.38 & \res[places=5]{build/R3d/faithful/standard/n5}{eta_ideal} & $\res[places=5]{build/R3d/faithful/standard/n5}{eta}\resflag{build/R3d/faithful/standard/n5}{trusted}{}{^{\dagger}}$ & \checkmark \\
        5a & R4a & 3 & 1.35 & \res[places=5]{build/R4a/faithful/standard/n3}{eta_ideal} & $\res[places=5]{build/R4a/faithful/standard/n3}{eta}\resflag{build/R4a/faithful/standard/n3}{trusted}{}{^{\dagger}}$ & \checkmark \\
        5b & R4b & 3 & 14.74 & \res[places=5]{build/R4b/faithful/standard/n3}{eta_ideal} & $\res[places=5]{build/R4b/faithful/standard/n3}{eta}\resflag{build/R4b/faithful/standard/n3}{trusted}{}{^{\dagger}}$ & \checkmark (one unit low) \\
        5c & R4c & 4 & 0.82 & \res[places=5]{build/R4c/faithful/standard/n4}{eta_ideal} & $\res[places=5]{build/R4c/faithful/standard/n4}{eta}\resflag{build/R4c/faithful/standard/n4}{trusted}{}{^{\dagger}}$ & \checkmark \\
        5d & R4d & 4 & 77.51 & \res[places=5]{build/R4d/faithful/standard/n4}{eta_ideal} & $\res[places=5]{build/R4d/faithful/standard/n4}{eta}\resflag{build/R4d/faithful/standard/n4}{trusted}{}{^{\dagger}}$ & \checkmark \\
        \midrule
        6b & R5a & 3 & 5.21559 & \res[places=5]{build/R5a/faithful/standard/n3x3}{eta_ideal} & $\res[places=5]{build/R5a/faithful/standard/n3x3}{eta}\resflag{build/R5a/faithful/standard/n3x3}{trusted}{}{^{\dagger}}$ & $\times$ \\
        6d & R5b & 4 & 603.863 & \res[places=5]{build/R5b/faithful/standard/n4x4}{eta_ideal} & $\res[places=5]{build/R5b/faithful/standard/n4x4}{eta}\resflag{build/R5b/faithful/standard/n4x4}{trusted}{}{^{\dagger}}$ & $\times$ \\
        6f & R5c & 5 & 216.469 & \res[places=5]{build/R5c/faithful/standard/n5x5}{eta_ideal} & $\res[places=5]{build/R5c/faithful/standard/n5x5}{eta}\resflag{build/R5c/faithful/standard/n5x5}{trusted}{}{^{\dagger}}$ & $\times$ \\
        7b & R6a & 2 & 3.75367 & \res[places=5]{build/R6a/faithful/standard/n2}{eta_ideal} & $\res[places=5]{build/R6a/faithful/standard/n2}{eta}\resflag{build/R6a/faithful/standard/n2}{trusted}{}{^{\dagger}}$ & $\times$ \\
        7c & R6b & 2 & 0.07984 & \res[places=5]{build/R6b/faithful/standard/n2}{eta_ideal} & $\res[places=5]{build/R6b/faithful/standard/n2}{eta}\resflag{build/R6b/faithful/standard/n2}{trusted}{}{^{\dagger}}$ & $\times$ \\
        \bottomrule
    \end{tabular}
\end{table}

\placeholderfigure{Printed choices distort the paper's solutions at the printed resolution.}{fig:faithful-vs-corrected}{Four panels, the paper's Figs 5a--5d (R4a--R4d): the analytic solution, the paper-faithful ground state (rounded constraint points, Maclaurin sources) and R4a--R4d solved with exact roots and Chebyshev sources, at the printed $n$; inset or second row: $L^\infty$ error against $n$ for both variants. Source: notebook \texttt{10\_paper\_reproduction}. Classical computation.}
% REVISE (T-27): this caption used to say "the corrected ground state" for the
%   exact-root/Chebyshev-source variant. ThesisPlanning.md's controlled
%   vocabulary retires "corrected" as a general term (only "R2a" vs "R2a as
%   printed" exists, confined to this chapter); the replacement above names
%   the variant by what changed, not by a corrected/faithful label.

\subsection{The Pipeline Is Faithful}
\label{sec:results-faithful}
% ~1.5 pp (S3). NEW in v2.0: merges the old "Operator Verification and
%   Encoding Cost" (verification half) and "State Preparation and Readout"
%   sections (ThesisPlanning.md §4.3 step 4). The encoding-COST half of the
%   old operators section has moved to sec:results-scaling (§8.3), since cost
%   is Q(i)/Q(ii) evidence, not a faithfulness check.
% LAND, in order:
%   1. Operators equal the assembled residuals in every mode
%      (fig:verification-accuracy, nb 20).
%   2. The circuit equals the emulation before normalising, <= 1e-10, on
%      every T1 record (nb 31).
%   3. Every trusted preparation within epsilon and its a priori bound
%      (nb 30); the decoded field within its budget (fig:decoded-field,
%      nb 32).
%   4. What the answer-independent choices cost: gamma for the initial
%      states, and an underestimated gap (nb 33).
%   5. The filter study in one table and two sentences (nb 34; B-17):
%      minimax needs 1.67-1.92x fewer degrees at equal suppression, in both
%      forms.
% FIG: fig:verification-accuracy (F7), fig:decoded-field (F8).

\placeholderfigure{Every encoding is verified to its mode's tolerance at every size tested.}{fig:verification-accuracy}{Verification error against $n$ for the published, structured and descriptor encodings: exact mode (filled markers), probe mode (open markers), IR mode (crosses), with each mode's tolerance as a dashed line; a second panel gives transpiled CNOT count per query against $n$ with fitted exponents. Source: notebook \texttt{20\_pillar1\_operators}.}

\placeholderfigure{The prepared state reproduces the analytic solution to within the predicted error budget.}{fig:decoded-field}{Decoded field $f_Q(x)$ against the analytic solution for R1 and R2a (both forms): circuit simulation (T1, with shot-sampled readout) and exact emulation (T2) overlaid; second panel: the error budget of \Cref{eq:error-budget} broken into $\varepsilon_{\mathrm{disc}}$, $\varepsilon_{\mathrm{filter}}$, $\varepsilon_{\mathrm{trunc}}$, $\varepsilon_{\mathrm{phase}}$ and $\varepsilon_{\mathrm{sample}}$ against $n$. Source: notebook \texttt{30\_pillar2\_stateprep}.}

\subsection{Where the Cost Lies}
\label{sec:results-scaling}
% ~2.0 pp -- THE HEADLINE FIGURE (F9); this single figure carries C and S2
%   (Q(i), Q(ii)). Label kept as sec:results-scaling (cited from Chs 4-5 and
%   the Abstract).
% LAND, in order:
%   1. THE HEADLINE TABLE (tab:headline): published | structured | descriptor,
%      over ||H||, the gap, alpha/||.||, cx per query, ancilla, exactness, the
%      degree at n = 8 as built and at alpha = ||R||, and the degree's
%      exponent. The Abstract references it. FABLE does not appear.
%   2. fig:scaling-headline (nb 40): against N, the encodings' looseness
%      alpha_H/||H||, the Hamiltonians' relative gap, and the degree (which
%      grows as the square root of the first over the second), for the
%      published, structured and descriptor encodings, with fitted exponents;
%      T2 points where measured and T3 beyond, VISIBLY distinguished.
%   3. Interpretation: the ENCODER WALL (published to structured) and the
%      HAMILTONIAN WALL (structured to descriptor) as two gaps on one plot.
%   3b. MOVED here from 5_descriptor.tex (2026-10-06): the measured growth of the degree between n = 7 and
%      n = 8, all T3 resource estimates. R2a: the descriptor's N^{res stateprep/R2a/descriptor/growth/T3/
%      filter=minimax degree} against the standard form's N^{res .../R2a/standard/growth/...}; R3a (leading
%      coefficient vanishes): descriptor .../R3a/descriptor/growth/...; R5b (two axes): descriptor
%      .../R5b/descriptor/growth/... against standard .../R5b/standard/growth/.... Two-point slopes of a
%      degree with log factors: say "consistent with", and fit more n if the budget allows. State plainly that
%      the descriptor's degree scales better on R2a (about N against N^4), and that this is the degree, not
%      the gates per query or the ancilla at reachable n (sec:descriptor-cost).
%   4. Resource fits against n with fitted exponents (checking claims B-03,
%      B-04, B-05, C-01, C-02 and A-07); the published build against the
%      paper's claims (A-06 to A-08): one verdict each.
%   5. "Could not simulate" is never "cannot reach".
% FIG: fig:scaling-headline (F9). TAB: tab:headline (T7).
% CHECK at results-v1: ThesisPlanning.md §3.9 labels the degree row "QITE
%   degree", but T-23 makes minimax primary; state which filter the final
%   numbers use (RESOLVED: P8's comparison is minimax).

\begin{table}[htbp]
    \centering
    \caption{Where the cost lies: the paper's regime under its published encoding and under the structured exact encoding, against the descriptor form, for Fig.~3b ($\{1,4,4\}$, $f(-1)=0$). Tier: T2 where measured, T3 beyond.}
    \label{tab:headline}
    \small
    \begin{tabular}{@{}l l l l@{}}
        \toprule
        Quantity & Standard, published & Standard, structured & Descriptor \\
        \midrule
        $\|H\|$ & \prov{$\Theta(N^8)$} & \prov{$\Theta(N^8)$} & \prov{$\Theta(N^2)$} \\
        Gap $\Delta$ & \prov{7.04 (flat)} & \prov{7.04 (flat)} & \prov{0.454 (flat)} \\
        $\alpha_{\G}/\|\G\|_2$ & \prov{$\sim N^3/2$} & \prov{$\le 2.12$} & no $\G$ \\
        $\alpha_H/\|H\|$ & \prov{$\sim N^{12}$} & \prov{22 $\to$ 60} & \prov{5.3 $\to$ 1.05} \\
        Gates per query & \prov{$O(n^3)$, approximate} & \prov{$O(n^2)$, exact} & \prov{$O(n^2)$, exact} \\
        Ancilla & \prov{$2n+3$ ($\G$ alone)} & \prov{$n + O(\log n)$ per $\G$} & \prov{$\lceil\log_2(n+4)\rceil + 8$} \\
        Degree at $n = 8$ ($\varepsilon = 10^{-6}$) & \prov{$8.2\times10^{23}$} & \prov{$1.2\times10^{10}$} & \prov{$5.6\times10^{3}$} \\
        Degree scaling & \prov{$\sim N^{10}$} & \prov{$\tilde\Theta(N^4)$} & \prov{$\tilde\Theta(N)$} \\
        \bottomrule
    \end{tabular}
\end{table}
% REVISE: the descriptor's "Ancilla" cell was + 6; CLAUDE.md Standing Facts
%   (Encoding costs) give ceil(log2(n+4)) + 8 for R2's system rows (already
%   corrected above) -- confirm against results-v1's actual record before
%   removing this comment.

\placeholderfigure[7cm]{The encoder wall and the Hamiltonian wall are separate gaps.}{fig:scaling-headline}{Three log--log panels against $N = 2^n$ for Fig.~3b: (a) $\|H\|$ and $\alpha_H$ for the published, structured and descriptor encodings; (b) the gap $\Delta$; (c) the filter degree $d$, with fitted exponents ($\approx N^{10}$, $\tilde\Theta(N^4)$, $\tilde\Theta(N)$). Exact emulation (T2) as filled markers, resource estimates (T3) as open markers. Source: notebook \texttt{40\_scaling\_and\_claims}.}

\placeholderfigure[7cm]{Three representative problems: where the comparison, quoted as built, reverses at the ideal encoding bound.}{fig:cost-reversals}{Case by case (R2a, R2c, R3a, R5b, ...), the descriptor's cost over the standard form's, as built and at the ideal bound $\alpha = \|R\|$, with the reversals marked (R3 at the ideal bound; the heat panel at small $n$). Optional. Source: notebook \texttt{40\_scaling\_and\_claims}.}

\subsection{The Descriptor's Gap and Its Rescaling}
\label{sec:results-gap}
% ~1.75 pp -- THE CHALLENGE (sec:descriptor-gap, sec:descriptor-rescaling made
%   a results section). NEW in v2.0, folding in the old "Representative
%   Problems" section's R2c material (ThesisPlanning.md §4.3 step 4).
% LAND, in order:
%   1. fig:gap-rescaling (nb 42): the gap and gamma^2 unscaled against
%      rescaled, every panel; the a priori rule against the best zeta on a
%      grid; where it stops (two axes, the singular ||H|| growth).
%   2. fig:gap-cases (nb 43): THREE representative problems in both forms,
%      each with its solution and cost -- R2c (stiff: rescaling's biggest
%      gain), R3a (a leading coefficient with zeros: where the ideal-bound
%      comparison reverses), R5b (heat on two axes: the gap still falls, and
%      the descriptor costs more at small n). The rest -> Appendix J
%      (app:catalogue), with a pointer naming what is there.
% FIG: fig:gap-rescaling (F11), fig:gap-cases (F12, reusing the old
%   fig:representative's three-problem layout).
% TODO: F11 is new in v2.0 (the gap/rescaling scaling figure did not exist
%   under the old "Representative Problems" framing); F12 can likely reuse
%   the old fig:representative's design (below) once the gap-specific framing
%   is drafted.

\placeholderfigure[7cm]{Block rescaling restores an O(1) gap on one axis, and stops at a double zero or on two axes.}{fig:gap-rescaling}{Panels across the benchmark ladder: the gap $\Delta$ and $\gamma^2$ unscaled against rescaled; the a priori $\zeta$ against the best $\zeta$ found on a grid; where the rescaling stops (a vanishing leading coefficient, a second axis). Source: notebook \texttt{42\_descriptor\_gap}.}

\placeholderfigure[7cm]{Three representative problems, solved in both forms.}{fig:gap-cases}{Three columns (R2c, R3a, R5b or R5c): top row, the decoded solution against the analytic one; bottom row, error against $n$ for the standard and descriptor forms (and, for R2c, with and without block rescaling). Exact emulation (T2). Source: notebook \texttt{40\_scaling\_and\_claims}.}

\subsection{Nonlinear Problems}
\label{sec:results-nonlinear}
% ~1.5 pp (Q(iii), S4).
% LAND, in order:
%   1. R6's kernel beside R8's (the projection evidence: kernel dimensions
%      against one) (nb 61).
%   2. R8: the K x n_t error surface and the degree comparison, both forms
%      (claim C-09).
%   3. Burgers (fig:nonlinear-burgers, nb 62): the lift's error per order
%      follows a/nu (about a/5nu) and NOT rho, which grows with the modes;
%      past a/nu ~ 5 the lift diverges while its kernel stays unique and
%      trusted: a NEGATIVE FINDING with its analysis.
% FIG: fig:nonlinear-burgers (F13, split from the old shared fig:carleman-ns).

\placeholderfigure[7cm]{Carleman truncation converges geometrically once time is resolved; the doubled space only projects; on Burgers the error tracks $a/\nu$, not $\rho$.}{fig:nonlinear-burgers}{(a) R8: error against Carleman order $K$ for several time resolutions $n_t$, both forms; (b) R6: projection fidelity against physical fidelity; (c) R9 Burgers: error per order against $a/\nu$, diverging past $a/\nu \approx 5$ while the kernel stays unique and trusted. Source: notebook \texttt{50\_navier\_stokes}.}
% REVISE: this figure used to be fig:carleman-ns, shared with
%   sec:results-ns. §5.8 points 5-6 give Burgers its own figure (F13) and
%   Navier--Stokes its own (F14, sec:results-ns below); split accordingly
%   once nb 62/nb 70 produce separate panels.

\subsection{Case Study: Two-Dimensional Navier--Stokes}
\label{sec:results-ns}
% ~1.25 pp.
% LAND, in order:
%   1. fig:navier-stokes (nb 70): vorticity snapshots and the degree pairs.
%      TAYLOR-GREEN validates F_1, the rows and the time axis (the lift exact
%      to rounding).
%   2. The vortex pair at a/nu = 3, 10 and 30 on 12-80 modes: the lift's
%      convergence and the truncation against a direct simulation; every
%      exact emulation within epsilon; the descriptor's time axis lowers the
%      degree in every pair, the ratio narrowing as the modes grow (C-10).
%   3. The TENSOR LAYOUT to 316 modes in 2K - 1 terms, with its gates per
%      query against the matchings (B-19).
%   4. T1 is BEYOND THE SIMULATION BUDGET here, with its estimate quoted:
%      said plainly.
% FIG: fig:navier-stokes (F14, split from the old shared fig:carleman-ns).

\placeholderfigure[7cm]{Navier--Stokes is reached by exact emulation and a resource estimate, not by circuit simulation.}{fig:navier-stokes}{(a) Taylor--Green and vortex-pair vorticity snapshots against DNS, with the tier marked (T2 small, T3 beyond); (b) the degree pairs, standard against descriptor time axis, narrowing as the modes grow; (c) the tensor-layout coupling's gates per query against the matchings, to 316 modes. Source: notebook \texttt{50\_navier\_stokes}.}

\subsection{Summary of Findings}
\label{sec:results-summary}
% ~0.5 pp. LAND: the seven numbered findings below (ThesisPlanning.md §5.8
%   point 7), AT LEAST TWO NEGATIVE (items 1, 3, 4 and 6 qualify). Then the
%   spine sentence (ThesisPlanning.md §3.3).
%   1. The paper's ODE results reproduce, its PDE and nonlinear ones do not.
%   2. The cost of the standard form is its Hamiltonian's.
%   3. The descriptor's degree advantage: large on one axis, modest on two,
%      REVERSED on the heat panel at small n and on R3 at the ideal bound.
%   4. The descriptor's gap erodes where derivatives are uncontrolled;
%      rescaling restores it on one axis only.
%   5. The pipeline is verified end to end at every size it reaches.
%   6. Carleman prepares in both forms, but a trusted kernel can hide a
%      diverging lift.
%   7. Navier--Stokes is reached by exact emulation and estimate, not by
%      circuit simulation.

% CHECKLIST (marker)
%   1. Is every number's tier marked?
%   2. Is "could not simulate" kept apart from "cannot reach"?
%   3. Are there at most three in-depth problems (sec:results-gap)?
%   4. Does every caption stand alone?
%   5. Are there at least two negative findings?
%   6. Is every Abstract number backed here?
%   7. Is sec:results-reproduction framed as a test of the paper, not rivalry?
